\section{Evaluation}

In this section, we evaluate the performance of GLM-OCR against current state-of-the-art pipeline tools, general Vision-Language Models (VLMs), and specialized OCR VLMs. To provide a comprehensive assessment, the evaluation is divided into two parts: standard Public Benchmarks and custom In-House Benchmarks.

\subsection{Public Benchmarks}

We first evaluate GLM-OCR on widely recognized public datasets encompassing Document Parsing and KIE tasks.

\paragraph{Overall Benchmark Performance.} As shown in Table~\ref{tab:docparse}, GLM-OCR demonstrates superior performance across the majority of standard datasets. In Document Parsing, our model achieves the highest scores on OmniDocBench v1.5 (94.6), OCRBench Text (94.0), UniMERNet (96.5), and TEDS\_TEST (86.0). It remains highly competitive on PubTabNet (85.2), trailing only MinerU 2.5. Furthermore, GLM-OCR establishes a clear SOTA in KIE, outperforming all available open-source competitors on Nanonets-KIE (93.7) and Handwritten-KIE (86.1), and even narrowing the gap with closed-source giants like Gemini-3-Pro.

% 插入表1
\begin{table*}[t]
\centering
\caption{Performance comparison of various OCR models. The results of Gemini-3-Pro and GPT-5.2-2025-12-11 are in gray as they are provided for reference only and excluded from the best-score ranking. The best results among the evaluated models are highlighted in bold.}
\label{tab:docparse}:
\resizebox{\textwidth}{!}{
\begin{tabular}{l ccccc cc}
\toprule
\textbf{Dataset} & \textbf{GLM-OCR} & \makecell{\textbf{PaddleOCR}\\\textbf{-VL-1.5}} & \makecell{\textbf{Deepseek}\\\textbf{-OCR2}} & \makecell{\textbf{MinerU}\\\textbf{2.5}} & \textbf{dots.ocr} & \makecell{\textcolor{gray}{\textbf{Gemini-3}}\\\textcolor{gray}{\textbf{-Pro}}} & \makecell{\textcolor{gray}{\textbf{GPT-5.2}}\\\textcolor{gray}{\textbf{-2025-12-11}}} \\
\midrule
\multicolumn{8}{l}{\textit{\textbf{Document Parsing}}} \\
\midrule
OmniDocBench v1.5 & \textbf{94.6} & 94.5 & 91.1 & 90.7 & 88.4 & \textcolor{gray}{90.3} & \textcolor{gray}{85.4} \\
OCRBench (Text)   & \textbf{94.0} & 75.3 & 34.7 & 75.3 & 92.1 & \textcolor{gray}{91.9} & \textcolor{gray}{83.7} \\
UniMERNet         & \textbf{96.5} & 96.1 & 85.8 & 96.4 & 90.0 & \textcolor{gray}{96.4} & \textcolor{gray}{90.5} \\
PubTabNet         & 85.2 & 84.6 & -    & \textbf{88.4} & 71.0 & \textcolor{gray}{91.4} & \textcolor{gray}{84.4} \\
TEDS\_TEST        & \textbf{86.0} & 83.3 & -    & 85.4 & 62.4 & \textcolor{gray}{81.8} & \textcolor{gray}{67.6} \\
\midrule
\multicolumn{8}{l}{\textit{\textbf{Key Information Extraction}}} \\
\midrule
Nanonets-KIE      & \textbf{93.7} & -    & -    & -    & -    & \textcolor{gray}{95.2} & \textcolor{gray}{87.5} \\
Handwritten-KIE   & \textbf{86.1} & -    & -    & -    & -    & \textcolor{gray}{94.5} & \textcolor{gray}{78.2} \\
\bottomrule
\end{tabular}
}
\end{table*}

\paragraph{OmniDocBench v1.5 Analysis.} To better understand the model's document parsing capabilities, we present a granular breakdown of the OmniDocBench v1.5 results in Table~\ref{tab:main_results}. This benchmark compares pipeline tools, general VLMs of varying sizes, and specialized VLMs.

Remarkably, despite possessing only 0.9B parameters, GLM-OCR achieves the highest Overall score (94.62), outperforming not only direct specialized competitors like PaddleOCR-VL-1.5 (94.50) and MinerU2.5 (90.67) but also massive general VLMs such as Qwen3-VL-235B (89.15) and Gemini-3 Pro (90.33).

A breakdown of sub-metrics indicates strong performance in table structure recovery. It achieves the absolute best scores in table recognition, scoring 93.96 on Table$_{TEDS}$ and 96.39 on Table$_{TEDS-S}$. While PaddleOCR-VL-1.5 slightly edges out GLM-OCR in Text$_{Edit}$ (0.035 vs 0.040) and Formula$_{CDM}$ (94.21 vs 93.90), GLM-OCR's exceptional table parsing capabilities secure its position as the top-performing model overall, proving that specialized, parameter-efficient architectures can rival or surpass scale-heavy models in complex document parsing.

% 插入表2
\begin{table*}[t]
\centering
\caption{Performance comparison with pipeline tools, general VLMs and specialized VLMs on OmniDocBench v1.5.}
\resizebox{\textwidth}{!}{
\begin{tabular}{l l c c c c c c c}
\toprule
\textbf{Model Type} & \textbf{Methods} & \textbf{Params} 
& \makecell{\textbf{Overall}$\uparrow$} 
& \makecell{\textbf{Text}$\downarrow$ \\ \textbf{Edit}}
& \makecell{\textbf{Formula}$\uparrow$ \\ \textbf{CDM}}
& \makecell{\textbf{Table}$\uparrow$ \\ \textbf{TEDS}} 
& \makecell{\textbf{Table}$\uparrow$ \\ \textbf{TEDS-S}} 
& \makecell{\textbf{Reading Order}$\downarrow$ \\ \textbf{Edit}} \\
\midrule

\multirow{3}{*}{Pipeline Tools}
& Marker-1.8.2~\cite{Marker-1.8.2} & - & 71.30 & 0.206 & 76.66 & 57.88 & 71.17 & 0.250 \\
& Mineru2-pipeline~\cite{Mineru2-pipeline} & - & 75.51 & 0.209 & 76.55 & 70.90 & 79.11 & 0.225 \\
& PP-StructureV3~\cite{PaddleOCR} & - & 86.73 & 0.073 & 85.79 & 81.68 & 89.48 & 0.073 \\
\midrule

\multirow{8}{*}{General VLMs}
& GPT-4o~\cite{GPT-4o} & - & 75.02 & 0.217 & 79.70 & 67.07 & 76.09 & 0.148 \\
& InternVL3-76B~\cite{zhu2025internvl3} & 76B & 80.33 & 0.131 & 83.42 & 70.64 & 77.74 & 0.113 \\
& InternVL3.5-241B~\cite{wang2025internvl35} & 241B & 82.67 & 0.142 & 87.23 & 75.00 & 81.28 & 0.125 \\
& GPT-5.2~\cite{gpt5_2} & - & 85.50 & 0.123 & 86.11 & 82.66 & 87.35 & 0.099 \\
& Qwen2.5-VL-72B~\cite{bai2025qwen2} & 72B & 87.02 & 0.094 & 88.27 & 82.15 & 86.22 & 0.102 \\
& Gemini-2.5 Pro~\cite{gemini25} & - & 88.03 & 0.075 & 85.82 & 85.71 & 90.29 & 0.097 \\
& Qwen3-VL~\cite{Qwen3-VL} & 235B & 89.15 & 0.069 & 88.14 & 86.21 & 90.55 & 0.068 \\
& Gemini-3 Pro~\cite{gemini30} & - & 90.33 & 0.065 & 89.18 & 88.28 & 90.29 & 0.071 \\
\midrule

\multirow{15}{*}{Specialized VLMs}
& Dolphin~\cite{feng2025dolphin} & 0.3B & 74.67 & 0.125 & 67.85 & 68.70 & 77.77 & 0.124 \\
& OCRFlux-3B~\cite{OCRFlux2025} & 3B & 74.82 & 0.193 & 68.03 & 75.75 & 80.23 & 0.202 \\
& Mistral OCR~\cite{mistral} & - & 78.83 & 0.164 & 82.84 & 70.03 & 78.04 & 0.144 \\
& POINTS-Reader~\cite{liu2025points} & 3B & 80.98 & 0.134 & 79.20 & 77.13 & 81.66 & 0.145 \\
& olmOCR-7B~\cite{poznanski2025olmocr} & 7B & 81.79 & 0.096 & 86.04 & 68.92 & 74.77 & 0.121 \\
& Dolphin-1.5~\cite{feng2025dolphin} & 0.3B & 83.21 & 0.092 & 80.78 & 78.06 & 84.10 & 0.080 \\
& MinerU2-VLM~\cite{Mineru2-pipeline} & 0.9B & 85.56 & 0.078 & 80.95 & 83.54 & 87.66 & 0.086 \\
& Nanonets-OCR-s~\cite{Nanonets-OCR-S} & 3B & 85.59 & 0.093 & 85.90 & 80.14 & 85.57 & 0.108 \\
& MonkeyOCR-pro-1.2B~\cite{li2025monkeyocr} & 1.9B & 86.96 & 0.084 & 85.02 & 84.24 & 89.02 & 0.130 \\
& Deepseek-OCR~\cite{wei2025deepseek} & 3B & 87.01 & 0.073 & 83.37 & 84.97 & 88.80 & 0.086 \\
& MonkeyOCR-3B~\cite{li2025monkeyocr} & 3.7B & 87.13 & 0.075 & 87.45 & 81.39 & 85.92 & 0.129 \\
& dots.ocr~\cite{dotsocr} & 3B & 88.41 & 0.048 & 83.22 & 86.78 & 90.62 & 0.053 \\
& MonkeyOCR-pro-3B~\cite{li2025monkeyocr} & 3.7B & 88.85 & 0.075 & 87.25 & 86.78 & 90.63 & 0.128 \\
& MinerU2.5~\cite{niu2025mineru2} & 1.2B & 90.67 & 0.047 & 88.46 & 88.22 & 92.38 & 0.044 \\
& PaddleOCR-VL~\cite{PaddleOCR-VL} & 0.9B & 92.86 & 0.035 & 91.22 & 90.89 & 94.76 & 0.043 \\
& PaddleOCR-VL-1.5~\cite{PaddleOCR-VL-1.5} & 0.9B & 94.50 & \textbf{0.035} & \textbf{94.21} & 92.76 & 95.79 & \textbf{0.042} \\
\midrule
& \textbf{GLM-OCR} & 0.9B & \textbf{94.62} & 0.040 & 93.90 & \textbf{93.96} & \textbf{96.39} & 0.044 \\
\bottomrule
\end{tabular}}
\label{tab:main_results}
\end{table*}

\subsection{In-House Benchmarks}

To assess the robustness of GLM-OCR in highly complex, real-world industrial scenarios, we conducted evaluations on a custom suite of in-house benchmarks. These tasks include Code Document parsing, Real-world Table extraction, Handwritten Text recognition, Multilingual Text processing, Seal Recognition, and Receipt KIE.

% 插入表3
\begin{table*}[t]
\centering
\caption{Performance comparison on custom real-world scenarios. The results of Gemini-3-Pro and GPT-5.2-2025-12-11 are in gray as they are provided for reference only and excluded from the best-score ranking. The best results among the evaluated models are highlighted in bold.}
\label{tab:self_eval}
\resizebox{\textwidth}{!}{
\begin{tabular}{l ccccc cc}
\toprule
\textbf{Task} & \textbf{GLM-OCR} & \makecell{\textbf{PaddleOCR}\\\textbf{-VL-1.5}} & \makecell{\textbf{Deepseek}\\\textbf{-OCR2}} & \makecell{\textbf{MinerU}\\\textbf{2.5}} & \textbf{dots.ocr} & \makecell{\textcolor{gray}{\textbf{Gemini-3}}\\\textcolor{gray}{\textbf{-Pro}}} & \makecell{\textcolor{gray}{\textbf{GPT-5.2}}\\\textcolor{gray}{\textbf{-2025-12-11}}} \\
\midrule
Code Document     & \textbf{84.7} & 75.8 & 82.1 & 82.9 & 80.8 & \textcolor{gray}{86.9} & \textcolor{gray}{84.4} \\
Real-world Table  & \textbf{91.5} & 86.1 & -    & 70.8 & 81.8 & \textcolor{gray}{90.6} & \textcolor{gray}{86.7} \\
Handwritten Text  & 87.0 & \textbf{87.4} & 73.8 & 54.2 & 71.7 & \textcolor{gray}{90.0} & \textcolor{gray}{78.0} \\
Multilingual Text & \textbf{69.3} & 54.8 & 56.1 & 27.8 & 65.1 & \textcolor{gray}{86.2} & \textcolor{gray}{70.1} \\
Seal Recognition  & \textbf{90.5} & 42.2 & 40.4 & -    & 63.0 & \textcolor{gray}{91.3} & \textcolor{gray}{58.8} \\
Receipt KIE       & \textbf{94.5} & -    & -    & -    & -    & \textcolor{gray}{97.3} & \textcolor{gray}{83.5} \\
\bottomrule
\end{tabular}
}
\end{table*}

As detailed in Table~\ref{tab:self_eval}, GLM-OCR achieves the highest score in five out of six evaluated categories among the compared open-weight models. Most notably, GLM-OCR demonstrates a notable margin in challenging, niche domains:

\begin{itemize}
    \item \textbf{Seal Recognition:} GLM-OCR achieves an exceptional score of 90.5, outperforming the next best open-weight model (dots.ocr at 63.0) by a massive margin and performing competitively with Gemini-3-Pro (91.3).
    \item \textbf{Multilingual Text:} The model excels in diverse linguistic contexts, scoring 69.3 compared to PaddleOCR-VL-1.5's 54.8.
    \item \textbf{Complex Formatting:} It leads in Code Document parsing (84.7) and Real-world Table extraction (91.5), proving its utility in highly structured and noisy environments.
\end{itemize}

While PaddleOCR-VL-1.5 holds a marginal lead in Handwritten Text (87.4 vs. GLM-OCR's 87.0), GLM-OCR remains highly effective. Crucially, in practical application tasks like Receipt KIE, GLM-OCR (94.5) easily surpasses proprietary models like GPT-5.2 (83.5). These in-house results validate that GLM-OCR is not merely optimizing for academic datasets but is highly capable of generalizing to the noisy, variable conditions of real-world OCR deployments.